\documentclass[11pt,a4paper]{article}
\usepackage{fontspec}
\usepackage{amsmath,amssymb,amsthm}
\usepackage{mathtools}
\usepackage{bm}
\usepackage{booktabs}
\usepackage{geometry}
\usepackage{enumitem}
\usepackage{xcolor}
\usepackage[pdfencoding=auto,psdextra]{hyperref}

\geometry{margin=2.5cm}

\newcommand{\E}{\mathbb{E}}
\newcommand{\R}{\mathbb{R}}
\newcommand{\KL}{\mathrm{KL}}
\newcommand{\N}{\mathcal{N}}
\newcommand{\sg}{\mathrm{sg}}
\newcommand{\Xspace}{\mathcal{X}}
\newcommand{\Za}{\mathcal{Z}_{A}}
\newcommand{\Zb}{\mathcal{Z}_{B}}
\newcommand{\Enc}{\mathcal{E}}
\newcommand{\Dec}{\mathcal{D}}
\newcommand{\den}{\mathrm{den}}
\newcommand{\Jac}{\mathbf{J}}
\newcommand{\norm}[1]{\left\lVert #1 \right\rVert}

\theoremstyle{definition}
\newtheorem{definition}{Definition}
\newtheorem{proposition}{Proposition}
\newtheorem{remark}{Remark}
\newtheorem{observation}{Observation}

\title{\textbf{Pixel-Space One-Step Denoiser Matching for Cross-VAE On-Policy Distillation}\\[0.3em]
\large Working Analysis (Draft)}
\author{Jayce-Ping}
\date{}

\begin{document}
\maketitle

\begin{abstract}
This working note formalizes a concrete cross-VAE distillation scheme: the student rolls out a
clean latent, decodes it to an image; that image is re-encoded by \emph{each} model's own VAE,
re-noised to a common level $t$, denoised in \emph{one step} by that model, and decoded back to
pixels; the pixel-space discrepancy between the teacher's and the student's reconstructions is the
training signal. We show this is a symmetric, fully pixel-anchored instance of the ``shared-space''
route (companion note \texttt{avoid\_inverse\_transport\_analysis.tex}, Route~B): it needs
\emph{no} latent transport map $\mathsf{T}$ or its inverse. We characterize what it distills
(one-step $x_0$-prediction matching across noise levels $\Leftrightarrow$ velocity/score matching
$\Leftrightarrow$ probability-flow matching, transported through the decoders), its gradient
structure (only adjoints/VJPs of forward maps), its behavior across $t$, and its irreducible
reconstruction floor. We also contrast it with DMD.
\end{abstract}

\section{Setup}
\label{sec:setup}

Two frozen VAEs give encoders/decoders $(\Enc_A,\Dec_A)$ on teacher latent space
$\Za\subseteq\R^{d_A}$ and $(\Enc_B,\Dec_B)$ on student latent space $\Zb\subseteq\R^{d_B}$, over a
shared pixel space $\Xspace$. The teacher velocity $v_A$ (frozen) and the student velocity
$v_\theta$ each induce a \emph{one-step clean-sample estimator} (denoiser). With a noising schedule
$z^t=\alpha_t z^0+\sigma_t\varepsilon$ ($\alpha_0=1,\sigma_0=0$, $\varepsilon\sim\N(0,I)$), the
denoiser of a model with velocity $v$ is the affine image of its velocity,
\begin{equation}
\den(z^t,t)\;\coloneqq\;\widehat{z^0}(z^t,t)
\;=\;\text{(schedule-dependent affine map of } v(z^t,t)\text{)},
\label{eq:den-def}
\end{equation}
i.e.\ the model's ``one step to clean'' prediction. Write $\den_A$ for the teacher and
$\den_\theta$ for the student. We assume $d_A\ge d_B$, but the scheme below does not use this.

\section{The Scheme}
\label{sec:scheme}

\begin{definition}[Pixel-space one-step denoiser matching]
\label{def:scheme}
Let the student roll out a clean latent $z_B^{0}\sim\rho_\theta$ and decode it to an image
\begin{equation}
x\;\coloneqq\;\Dec_B\big(z_B^{0}\big)\ \in\ \Xspace \qquad(\text{on-policy image}).
\label{eq:img}
\end{equation}
For a noise level $t$ and independent draws $\varepsilon_A\in\R^{d_A},\ \varepsilon_B\in\R^{d_B}$,
define the two \emph{decode-of-one-step-denoise} branches
\begin{align}
\text{teacher:}\quad
y_A(x,t,\varepsilon_A)&\;\coloneqq\;\Dec_A\!\Big(\den_A\big(\alpha_t\,\Enc_A(x)+\sigma_t\varepsilon_A,\ t\big)\Big),
\label{eq:yA}\\
\text{student:}\quad
y_B(x,t,\varepsilon_B)&\;\coloneqq\;\Dec_B\!\Big(\den_\theta\big(\alpha_t\,\Enc_B(x)+\sigma_t\varepsilon_B,\ t\big)\Big).
\label{eq:yB}
\end{align}
The training objective is the pixel-space discrepancy
\begin{equation}
L(\theta)\;=\;\E_{x\sim\rho_\theta}\ \E_{t\sim p(t)}\ \E_{\varepsilon_A,\varepsilon_B}\Big[\,w(t)\,
\norm{\ \sg\!\big[y_A(x,t,\varepsilon_A)\big]-y_B(x,t,\varepsilon_B)\ }^2\Big],
\label{eq:loss}
\end{equation}
with a time weight $w(t)$ and a stop-gradient $\sg[\cdot]$ on the teacher target.
\end{definition}

\begin{observation}[No transport map is needed; pixels are the shared anchor]
\label{obs:no-transport}
Both branches only ever apply \emph{forward} maps within a single latent space
($\Enc_\bullet$, $\den_\bullet$, $\Dec_\bullet$), and meet in the shared pixel space $\Xspace$.
The two latent spaces are never put into correspondence: there is no $\mathsf{T}:\Za\to\Zb$ and,
in particular, no underdetermined inverse $\mathsf{T}^{-1}:\Zb\to\Za$. The cross-VAE mismatch is
dissolved by comparing \emph{behavior in pixel space} rather than \emph{coordinates in latent
space}. This is the symmetric, both-sides-decoded form of Route~B of the companion note.
\end{observation}

\section{What the Objective Distills}
\label{sec:what}

\begin{proposition}[Denoiser family $\Leftrightarrow$ velocity $\Leftrightarrow$ flow]
\label{prop:equiv}
For a fixed schedule, the one-step denoiser family $\{\den(\cdot,t)\}_{t\in(0,1)}$ is in bijection
with the velocity family $\{v(\cdot,t)\}_{t\in(0,1)}$ via the affine relation \eqref{eq:den-def};
the velocity family determines the probability-flow ODE and hence the model's clean marginal
$p^{0}$. Consequently, matching one-step $x_0$-predictions \emph{at all $t$} (within a common
space) is equivalent to matching velocities/scores, i.e.\ to matching the generated distribution.
\end{proposition}

\begin{remark}[The pixel version = velocity matching transported through the decoders]
\label{rem:transport-through-decoder}
Objective \eqref{eq:loss} compares $\Dec_A(\den_A(\cdot))$ against $\Dec_B(\den_\theta(\cdot))$.
Were the decoders invertible and identical, this would reduce (by Prop.~\ref{prop:equiv}) to exact
latent denoiser/velocity matching. In general it is velocity matching \emph{conjugated by the
decoders}: the comparison is transported to pixels by the \emph{forward} maps $\Dec_A,\Dec_B$
instead of by a latent transport map. This is precisely why the ill-posed vector-field pushforward
(and its section/inverse) is avoided---see companion note, Prop.~``projectability''.
\end{remark}

\begin{proposition}[Fixed point]
\label{prop:fixed-point}
Suppose $\Dec_B$ is injective on the relevant support and $w(t)>0$. If $L(\theta)$ attains its
$\theta$-dependent minimum with $y_B\equiv y_A$ (in expectation over $\varepsilon$) for all
$x\in\mathrm{supp}(\rho_\theta)$ and all $t$, then $\den_\theta$ reproduces the pixel-conjugated
teacher denoiser on that support, hence (Prop.~\ref{prop:equiv}) the student's probability flow
matches the teacher's \emph{as observed through} $\Dec_B$. As $\rho_\theta$ then covers the
teacher's image manifold, the match is self-reinforcing---the standard on-policy fixed point.
\end{proposition}

\section{Gradient Structure and Stop-Gradient Design}
\label{sec:grad}

The student parameter $\theta$ enters \eqref{eq:loss} through (i) the on-policy image $x$
(produced by the rollout $v_\theta$) and (ii) the student branch $y_B$ (the denoiser being
trained). The teacher branch $y_A$ depends on $\theta$ only through $x$.

\begin{remark}[Recommended gradient paths]
\label{rem:paths}
Treat $x$ as a \emph{sampled on-policy state} (detach it) and stop-gradient the teacher target;
then the gradient flows only through the student branch,
\begin{equation}
\nabla_\theta L
= \E\Big[-2\,w(t)\,\big(\sg[y_A]-y_B\big)^{\!\top}\,
\Jac_{\Dec_B}\,\Jac_{\den_\theta}\,\partial_\theta\den_\theta\Big],
\label{eq:grad}
\end{equation}
which uses only \emph{adjoints} (vector--Jacobian products) of the forward maps $\Dec_B$ and
$\den_\theta$; no map is inverted. Crucially, the rollout policy and the denoiser are the
\emph{same} network $v_\theta$, so regressing the denoiser on on-policy states \eqref{eq:grad}
\emph{also} improves the rollout---this is the on-policy distillation mechanism. Propagating
additionally through $x$ (the rollout) is possible but expensive and can create a trivial
fixed point (the student may prefer images on which the two VAEs happen to agree); detaching $x$
avoids this.
\end{remark}

\section{Behavior Across Noise Levels and the Reconstruction Floor}
\label{sec:t-behavior}

\begin{remark}[Limiting behavior in $t$]
\label{rem:t-limits}
At $t\to0$: $\den_\bullet\to\mathrm{id}$, so $y_A\to\Dec_A\Enc_A(x)$ and $y_B\to\Dec_B\Enc_B(x)$;
the loss tends to the VAE round-trip discrepancy, largely $\theta$-independent and
\emph{uninformative}. At $t\to1$: the branches compare the two models' \emph{generative priors}
(one-step guesses from near-noise); the distillation signal is strongest but has the highest
variance. A weight $w(t)$ emphasizing mid-to-high $t$ is therefore natural.
\end{remark}

\begin{proposition}[Irreducible floor from VAE mismatch]
\label{prop:floor}
Even a student whose implied clean distribution equals the teacher's incurs a nonzero loss,
because $y_A$ and $y_B$ pass through \emph{different} decoders. Define the reconstruction floor
\begin{equation}
\delta_{\mathrm{rec}}\;\coloneqq\;\E_{x\sim\rho_\theta}\ \norm{\Dec_A\Enc_A(x)-\Dec_B\Enc_B(x)},
\label{eq:floor}
\end{equation}
the pixel gap between the two VAEs' reconstructions of the same image. Then \eqref{eq:loss} is
bounded below by a $\delta_{\mathrm{rec}}$-type term (exactly $\delta_{\mathrm{rec}}^2$ in the
$t\to0$ limit) and should be interpreted \emph{modulo} this floor; e.g.\ one may subtract the
student's own reconstruction baseline to debias.
\end{proposition}

\section{Relationship to DMD, and to the Companion Routes}
\label{sec:relations}

\begin{remark}[DMD comparison]
\label{rem:dmd}
Like DMD/VSD, the states probed are \emph{noised outputs of the student generator} (here: noised
re-encodings of $x\sim\rho_\theta$), so the scheme is on-policy w.r.t.\ the student's \emph{output}
distribution (not its trajectory states). Unlike DMD, the signal is a \emph{direct pixel-space
regression} between decoded one-step denoisings \eqref{eq:loss}, back-propagated through the
student denoiser; it needs \emph{no} auxiliary online ``fake-score'' network and does not form the
$(\,\text{real score}-\text{fake score}\,)$ gradient. It is thus closer to a denoiser/$x_0$-matching
distillation than to variational score distillation, at the cost of back-propagating through
$\Dec_B$ each step.
\end{remark}

\begin{remark}[Placement among the routes]
\label{rem:placement}
This scheme is Route~B (pixel adjoint) made symmetric and swept over $t$: both sides do
encode$\to$noise$\to$one-step$\to$decode, and the target uses only the \emph{frozen}
$\Enc_A,\Dec_A$ (never a learned low$\to$high map). It preserves the student as a
velocity/denoiser network $v_\theta$ (a multi-step flow), in contrast to Route~A (DMD), which
naturally yields a few-step generator.
\end{remark}

\section{Assumptions and Limitations}
\label{sec:limits}

\begin{itemize}[leftmargin=1.4em,itemsep=2pt]
  \item \textbf{Metric.} The pixel-space $\ell_2$ metric, pulled back through $\Dec_B$, induces a
        non-Euclidean effective metric on $\Zb$ (weighted by $\Jac_{\Dec_B}^{\!\top}\Jac_{\Dec_B}$);
        perceptual/latent structure is only implicitly respected.
  \item \textbf{Floor.} $\delta_{\mathrm{rec}}$ \eqref{eq:floor} sets a noise floor; small-$t$ terms
        are dominated by it and should be down-weighted or debiased (Prop.~\ref{prop:floor}).
  \item \textbf{Variance.} $\varepsilon_A,\varepsilon_B$ are independent (dimensions differ), adding
        variance; averaging over several draws, or using expected (denoiser) forms, reduces it.
  \item \textbf{Cost.} Each step needs a teacher forward plus $\Enc_A,\Dec_A$ for the target and a
        differentiable pass through $\Enc_B,\den_\theta,\Dec_B$ for the student branch.
  \item \textbf{On-policy notion.} On-policy is w.r.t.\ the student \emph{output} distribution
        (noised), not the multi-step \emph{trajectory} states of XOPD's L1.
  \item \textbf{One-step coarseness.} At high $t$ the one-step $x_0$-prediction is a coarse
        estimate for both models; the signal is meaningful only when integrated across $t$
        (Prop.~\ref{prop:equiv}).
\end{itemize}

\section{Summary}
\label{sec:summary}

The scheme distills the teacher's denoiser (equivalently its velocity/flow) into the student by
comparing \emph{decoded one-step denoisings in pixel space}, on the student's own on-policy image
distribution, across noise levels. Its defining virtue is that pixel space serves as the shared
anchor, so it requires \emph{no} latent transport map and \emph{no} underdetermined
student$\to$teacher inverse; its gradient uses only adjoints of forward maps. Its principal
caveats are the VAE-mismatch floor $\delta_{\mathrm{rec}}$, the pixel-metric bias, and the per-step
decode cost. It keeps the student a multi-step flow model (unlike DMD) and needs no auxiliary
score network (unlike VSD).

\end{document}
